\chapter{From bank records to a reproducible study}
\label{r3:construction}

An economist studying a bank policy needs to reconstruct what customers
and decision makers knew at the time. Chapter~\ref{r3:product} defines the
comparison; the data must preserve it through reporting delays, corrections
and changes in customer identity. The same population, action and outcome
should be analysed whether the study runs in a research notebook, a scheduled
command or an analytical application. The construction argument proceeds from revision-aware histories to statistical reductions and proposed workspaces. Complete-workload measurements then test the costs of partitioning and recovery, before the language-specific protocol comparison examines which result semantics survive an application boundary.

\section{Constructing a history from economic and information time}
\label{r3:history-construction}

Represent an event revision by
$e=(k,j,v,t_e,t_a,x,u,z)$: stable event key $k$, entity $j$, revision $v$,
economic time $t_e$, availability time $t_a$, value $x$, unit $u$ and
void marker $z$. A decision has a stable key and time $\tau$.
Under an explicitly chosen inclusive availability convention, its known
revision of event $k$ is
\[
 v^*(k,\tau)=\max\{v:t_a(k,v)\le\tau\}.
\]
If the source has a different revision order, that order must replace
the numerical maximum. This construction admits availability at the decision cutoff only under that explicit inclusive convention. The strict feature-clock service discussed in Chapter~\ref{in:operations} excludes equality; exchanging the views therefore requires a declared clock mapping, not an assumption that their accepted rows coincide. For a lookback $\ell$, a monetary feature is
\begin{equation}
 X_j(\tau)=
 \sum_{\substack{k:\operatorname{entity}(k)=j\\v^*(k,\tau)\ {\rm exists}}}
 x_{k,v^*}\,
 \ind\{\tau-\ell<t_{e,k,v^*}\le\tau,\ z_{k,v^*}=0\}.
 \label{r3:revision-feature}
\end{equation}
Revision selection precedes economic-window filtering. If a later
correction moves an event beyond $\tau$, filtering time first would
resurrect an obsolete revision. Likewise a late-arriving past event
belongs to a retrospective accounting view but may be absent from the
history available to the historical policy.

The feature computation needs a source-completeness declaration. Zero
observed transactions can mean a complete feed with no activity, an
unconnected account or a failed batch. Only the first supports a zero
activity value under the stated feed contract. Missingness indicators,
known outages and event-count denominators should survive construction.
Treating each absent row as zero would silently replace the observation
operator in Chapter~\ref{r3:observation}.

A bank may have several decision clocks: quote, approval, acceptance,
settlement and account posting. Features available at approval need not
have been available at quotation. A timestamp equality may also conceal
unknown within-batch ordering. The initial convention can reject such
ambiguity or use a documented cutoff; it should not invent millisecond
order. Period-end stocks, period sums and period averages need separate
aggregation definitions, as the aliasing calculations demonstrated.

Entity linkage introduces another layer. Legal customer, household,
account, facility, branch and corporate group are distinct units. A
household action can affect several accounts; a branch rollout induces
shared assignment. Crosswalks need effective and availability dates.
When entity membership is uncertain, the product should preserve the
linkage hypothesis or weighted alternatives needed by the declared
analysis. A hard merge changes both potential outcomes and the
independent sampling units. The total-variance decomposition in
Section~\ref{r3:entity-uncertainty} shows why exporting one identifier
without its uncertainty can understate error.

\section{What automatic organization can and cannot infer}
\label{r3:automatic-organization}

Automatic profiling can propose event keys, candidate clocks, repeated
decisions, missingness patterns, currencies, ownership crosswalks and
aggregation rules. It can compare the proposal with physical data,
report violations and display the rows that justify a mapping.
An analyst then supplies the economic semantics: which field is the
action, when the bank knew a value, what the outcome means, and which
changes were possible. This produces a reusable observation specification
rather than an unexplained cleaned table.

The proposal process should use competing hypotheses when evidence is
ambiguous. For example, a table with several rows per customer-month may
contain revisions, different products, or repeated actions. Deduplicating
it under the first hypothesis can delete genuine decisions under the
third. The system can rank the hypotheses by key consistency and source
metadata, but causal conclusions require the selected interpretation.
A useful user view shows the consequences: population size, affected
histories, assignment timing and target units under each mapping.

Schema tools, columnar engines, feature stores and catalogs provide
parts of this workflow \citep{sw_duckdb,sw_polars,sw_pandera,sw_feast,
sw_datahub,sw_openmetadata,sw_splink}. Their source-specific limitations
were established in Section~\ref{r3:software-data}. The common product
should preserve the bank's existing asset identifiers and attach the
causal observation semantics. It should avoid maintaining a competing
catalog or treating catalog lineage as a discovered economic graph.

The output is a versioned trajectory or panel view with an explicit
loss-of-information declaration. For example, a monthly snapshot may
discard within-month action order. That view is suitable for a model
assuming one monthly decision, but cannot support a finer sequential
intervention merely by adding a lag column. The original event references
must remain recoverable so an alternative clock can be reconstructed.

\section{A scientific plan and a physical plan}
\label{r3:execution-plan}

The scientific plan fixes the target, model restrictions, observation
view, identification expression, estimator, inference construction,
statistical split and numerical tolerances. The physical plan chooses
data partitions, workers, memory limits, caching and reduction order.
Changes to the latter are permitted only when they preserve the former.
This is the execution counterpart of Equation~\eqref{r3:gradient-pullback}.

Each materialized object should identify its inputs and semantic version.
A result refers to the frozen scientific plan, data snapshot, qualified
backend revision, random seeds, numerical diagnostics and source lineage.
An input correction creates a new result lineage; it does not rewrite the
earlier answer. Deterministic content hashes are useful identities, but
they do not establish that two differently serialized scientific models
are equivalent. Canonicalization rules need their own specification.

For cross-fitting, a fold is a function of the permitted sampling unit,
assigned before scheduling. A household remains in the same fold when
its account rows move between workers. Multiway or blocked-time designs
need their own split objects. Hyperparameter selection, feature screening
and representation learning must follow the same training/evaluation
boundary as the nuisance model. Storing only a model checkpoint while
losing its training population makes later leakage review impossible.

Retries need contribution identities. A task can be retried without
adding its score twice if the reducer commits at most one completed
artifact for each tuple of analysis, fold, partition and attempt lineage.
This is an engineering idempotency requirement; a mathematical sum
alone does not supply it. Partial failure should leave a resumable plan
and an incomplete result, not a silently smaller study.

\section{The algebra that can be distributed}
\label{r3:distributed-algebra}

For precomputed vector scores $\psi_i\in\mathbb R^d$, a partition can
return counts, sums and derivative sums:
\[
 n_b,\quad U_b=\sum_{i\in b}\psi_i,\quad
 J_b=\sum_{i\in b}\partial_\theta\psi_i.
\]
Their sums reconstruct the global moment and Jacobian. For a cluster
covariance, it must additionally return partial sums keyed by the
statistical cluster. Regrouping those contributions yields
$U_c=\sum_{i\in c}\psi_i$. Outer products $U_cU_c'$ are formed after
that regrouping. Their order is essential, as the counterexample in
Section~\ref{r3:software-execution} shows.

When the contribution is an uncentered influence-value candidate $D_i$
for a mean $\psi$, retain $S_c=\sum_{i\in c}D_i$ and $n_c$.
The centered cluster sum is
$S_c-n_c\widehat\psi$, with $\widehat\psi=\sum_cS_c/\sum_cn_c$.
Centering each partition by its own mean destroys this identity.
For entity/date clustering, retain sums for entities, dates and observed
intersections. The inclusion--exclusion covariance can be indefinite
in a finite sample; silently projecting it to a positive matrix changes
the reported estimator. A qualified remedy requires an explicit rule
and justification.

Different estimators require different reductions. Linear IV/GMM can
accumulate moment cross-products, but nonlinear GMM may require repeated
passes and candidate-specific moments. A weak-IV inversion evaluates
a test over parameter values; replacing it by one inverse Hessian changes
the inference method. An empirical quantile, a survival curve and a
regularized inverse problem do not generally reduce to the same score
mean. The operation catalogue must advertise its actual sufficient
computational outputs.

For $C$ observed clusters and score dimension $d$, keyed sum storage is
$O(Cd)$ before matrix construction, and outer products cost $O(Cd^2)$.
Dense storage for $E$ entities by $D$ dates uses $O(EDd)$ even when few
intersections occur. A sparse implementation can reduce storage to
the observed intersections, but its hash/sort costs and data skew must
be measured. Streaming over rows is not a guarantee of bounded memory
as the number of statistical groups grows.

\section{State-space and discovery workloads scale differently}
\label{r3:dynamic-scale}

For $N$ independent trajectories, horizon $T$ and $K$ hidden states,
dense forward/backward evaluation costs $O(NTK^2)$ per likelihood pass.
An EM iteration can distribute trajectories and combine expected
transition/emission counts under a common parameter vector. Mixing counts
from different parameter iterations breaks the EM objective and its
ascent argument. Online or asynchronous updates are different algorithms
and need their own convergence conditions.

Within a single trajectory, messages have a chronological dependency.
Unnormalized transition/emission matrices can be composed associatively,
so a parallel prefix construction is possible, but dense matrix products
cost $O(K^3)$ instead of the $O(K^2)$ vector update and require careful
scaling. It is a performance tradeoff, not free parallelism. Smoothing
requires corresponding backward information. For continuous states,
particle count, state dimension, degeneracy and resampling communication
can dominate costs; likelihood error must be assessed in the downstream
parameter or policy units.

A learned structural model also has repeated optimizer, posterior or
regularization loops. Reporting only the speed of one filter pass
omits much of its cost. The relevant workload includes initialization,
multiple starts, rank diagnostics, parameter uncertainty and policy
evaluation. The approximation should be named: sparse transition,
low-rank factor, variational family, particle distribution or finite
state truncation. Each may change a scientific model or numerical error.

For $p$ observed series and maximum lag $L$, even storing candidate
lagged pairs requires order $p^2L$ objects. Conditional-independence
searches add conditioning-set complexity; dense matrix-exponential
acyclicity constraints have another scaling profile. Screening or
hierarchical grouping can make the workload feasible, but may remove
an adjustment path or a weak confounder. Any screening restriction
therefore belongs in the discovery model or its recovery argument.
Changing to monthly frequency is an observation change, not merely
a compute optimization.

The initial deployment should support local columnar execution for
bounded workloads, then qualified distribution where the algebra allows
it. Ray is a potential scheduling component \citep{sw_ray}; it does not
provide the cluster or mixing assumptions. GPU acceleration is most
plausible for dense numerical kernels and learned nuisance models.
Sorting temporal revisions, irregular groups and small exact
identification problems can have different bottlenecks. A hardware
choice should follow a measured workload and error tolerance.

\section{One core library and three standalone workspaces}
\label{r3:workspaces}

The standalone product should offer three initial workspaces sharing
the same request and evidence objects.

The \emph{intervention study} workspace starts with a population and a
proposed or historical decision. It constructs an encouragement, adjusted
effect or rollout analysis; shows overlap, timing and comparison groups;
and produces effects in outcome and financial units. A reviewer can
inspect which assumptions distinguish receipt from invitation and which
aggregation defines the reported business effect.

The \emph{sequential decision} workspace reconstructs customer histories,
compares observed-history policies and controlled-state explanations,
and displays feasible actions, support and value uncertainty. It
includes policy alternatives and the cost account. An attractive latent
policy cannot bypass its logger and identification review.

The \emph{dynamic systems} workspace combines panel and macro time
series, observation-frequency analysis, structural-shock alternatives
and discovery. It compares graph/model classes and their implied
response sets. A user can see why finer data, a new instrument or an
experiment would change the conclusion.

In each workspace the primary view should answer the user's decision:
what action is compared, for whom, over which horizon, at what expected
value and with which unresolved assumptions. Technical traces, scores
and solver logs remain available to researchers and validators without
becoming mandatory steps for every business reader. An embedded client
uses the same operations through a library or job service and can render
its own application-specific view.

\section{Interfaces for neighboring financial applications}
\label{r3:integration-contract}

Table~\ref{r3:integration} assigns responsibilities for proposed
integrations with the named research applications. These are interface
proposals; no working connection or acceptance by a sibling project
is assumed. The host remains responsible for its commercial model and
user decision, while the causal core preserves the requested estimand
and scientific evidence.

\begin{longtable}{P{.16\textwidth}P{.37\textwidth}P{.40\textwidth}}
\caption{Proposed reusable interfaces, with explicit ownership}
\label{r3:integration}\\
\toprule Host & Supplied to the causal core & Returned and retained by the host\\
\midrule\endfirsthead
\toprule Host & Supplied to the causal core & Returned and retained by the host\\
\midrule\endhead
FinMarketing &
Eligible population, invitation/receipt definitions, available features,
assignment history and product cash-flow account. &
Targeted effects or policy values, population/weight definitions,
uncertainty and unresolved assumptions. The host owns offer construction,
customer constraints and its product economics.\\
dataml &
Data snapshot references, candidate schema/clock/linkage interpretations
and observation uncertainty. &
Required information sets, consequences of alternative mappings,
accepted trajectory views and unresolved observation assumptions.
Neither system silently certifies ignorability.\\
BayesFilter &
The causal layer supplies a declared state likelihood, actions,
observation masks and numerical tolerances to a proposed numerical adapter. &
Likelihood/filter/posterior artifacts and approximation diagnostics.
The causal layer owns logger replacement and value identification;
numerical health does not establish those properties.\\
MacroFinance &
Series and release vintages, structural/shock definitions, model
restrictions, instrument normalization and observation frequency. &
Qualified moment/response analyses and identified alternatives.
The host owns its equilibrium and pricing restrictions; a predictive
innovation is not automatically an economic shock.\\
Other applications &
A versioned request using the same target, observation and sampling types;
local adapters for their data and outcomes. &
Portable result objects and reproducible references, without importing
another application's private business assumptions.\\
\bottomrule
\end{longtable}

The minimum integration exercise should round-trip a request, execute a
small supported study and reproduce its result independently. It should
then deliberately alter a currency, weight role, clock or graph endpoint
and demonstrate either a correctly changed target or a rejected conversion.
Only after that exercise should an adapter be called compatible.
Serializing the same array shape is insufficient.

Python can serve as the initial scientific orchestration layer while
data artifacts use language-neutral tables and structured metadata.
R or Julia backends can run in isolated, versioned jobs when their
qualified method is needed. A deployment must preserve package and
license boundaries identified in Chapter~\ref{r3:software}; some
components require an explicit license or service-design decision.
The core request/result representation should remain usable when a
backend is replaced or unavailable.

\section{Construction sequence and measurable acceptance}
\label{r3:construction-sequence}

The first construction stage implements immutable observation views,
scientific requests, identification traces and bounded result types.
It needs exact positive and negative fixtures before a general interface
is advertised. Stage two qualifies adjusted/encouragement/cohort methods
and finite-horizon observed-history values against independent numerical
references, including their cluster and policy-specific scores.

Stage three connects controlled-state fitting and discovery to the same
scientific objects. It starts with finite models, supported graph classes
and explicit unresolved outputs. Latent causal identification, proxy
stability, uncertainty propagation and structural choice remain visible
subtasks. They belong in the initial research architecture, even where
their numerical qualification takes longer than ordinary adjustment.

Stage four measures the complete workflow and tests reuse in the three
workspaces. For a fixed scientific request, measure data reconstruction,
nuisance/latent fitting, inference, sensitivity, review and report
production separately. Record wall time, peak memory, processor/accelerator
use, retries, spill volume and numerical disagreement against the
qualified reference. Throughput without the required diagnostics is a
different output and an invalid economic comparison.

Acceptance should require successful replay, deliberate failure behavior,
target-preserving adapter tests and appropriate repeated-sample checks.
Performance targets must be attached to a workload: number of entities,
decisions, variables, lags, clusters and latent states, rather than a
single bank-wide row count. Availability, access control, locality and
operational recovery are deployment requirements to be specified with a
particular bank; the proposed scientific core does not establish their
implementation.

This sequence permits a useful independent product and shared library
without assuming that all host applications use the same causal model.
Its central reusable asset is the preserved scientific meaning of an
analysis, supported by qualified computation and a traceable financial
decision. The following workload studies test selected parts of this construction. Their serial, partitioned and restartable comparisons retain different timing boundaries and statistical laws; they are complementary evidence, not interchangeable speed estimates.
\section{Computational scale follows the statistical construction}
\label{s5:dominant-workloads}
Ten million precomputed scores say little about the cost of fitting a latent
decision model. A scientific workload must retain the dimensions that make
its inference difficult. For responses these include dates, horizons and
covariance coordinates; for controlled states, the number of states,
trajectory length, validation sample and complete refits; for discovery,
variables, lag window and conditioning-set search. The expanded measurements
therefore execute the operations developed in the preceding chapters.

\subsection{Measured work and numerical outputs}
Each case runs twice in a fresh process on one AMD EPYC 7773X host, with one
thread in every observed numerical pool. Synthetic repetitions use separate
fixed seeds; the public case changes its declared fold seed. The ranges
therefore include data-dependent algorithm cost and host variation, not just
identical-input timing noise. Process elapsed time includes interpreter and
module loading, generation or public-data reading, local write/read,
estimation, the specified uncertainty calculation and report serialization.
It excludes analyst review, remote storage, bank access controls and a
distributed scheduler. Memory is process peak resident memory.

\begin{center}\small
\begin{tabular}{P{.57\textwidth}rr}\toprule
Scientific workload & Seconds & Peak MiB\\\midrule
Responses: 4,000 dates, horizon 12 & .77--.78 & 100\\
Responses: 16,000 dates, horizon 36 & 1.19--1.22 & 153--154\\
Private finite states: $K=3,N=300,H=8$ & 2.12--2.26 & 74--76\\
Private finite states: $K=4,N=600,H=12$ & 7.49--7.64 & 78--79\\
Private finite states: $K=6,N=3,000,H=24$ & 121.66--132.50 & 151--153\\
Temporal search: 6 variables, 2 lags, 1,600 dates & 1.34--1.44 & 104\\
Temporal search: 12 variables, 4 lags, 3,200 dates & 4.28--5.23 & 106\\
Temporal search: 24 variables, 8 lags, 6,400 dates & 50.59--101.89 & 128--158\\
Observed history: 2,400 customers, 4 decisions & 2.37--2.38 & 188--192\\
Gaussian controlled state: 800 customers, 12 decisions & .93--.95 & 96\\
Public finance: 9,915 households, ATE and LATE & 3.83--3.84 & 283--286\\\bottomrule
\end{tabular}\end{center}
Responses compute all three covariance alternatives and weak-denominator
sets. Each finite-state run includes a two-start joint fit, observed
information and three refits of whole trajectories and validation data. The three draws
measure cost; they do not estimate confidence-interval tails. The state grids
use respectively four, five and seven signals and 300, 600 and 1,500
validation measurements. All six main fits and eighteen refits converge
without a reported parameter boundary; main information calculations are
available. This numerical availability does not supersede the undercoverage
experiment in Section~\ref{s5:private-logger}.

In the largest finite model, fitting and information take 45.94--48.55
seconds, while three refits take 75.28--83.49 seconds. In the largest
discovery model, PCMCI+ takes 7.74--8.05 seconds and LPCMCI takes
41.98--92.97 seconds, with the conditioning cap of three retained.
These searches preserve graph and test matrices; their output is an
exploratory graph, not an automatically qualified effect interval. The
history workload includes independent policy training, backwards nuisance
fitting and contrast covariance. The Gaussian workload includes the joint
logger likelihood, analytic score, information and policy value.

The public workflow fits eligibility and participation nuisances, constructs
the ratio set, and checks DoubleML using the same predictions. The matched
package call adds about .043 seconds after those shared fits; it excludes a
second nuisance-learning pass. It is therefore a semantic cross-check, not
evidence that the product is faster than an expert using DoubleML directly.
Installed versions are Python 3.13.13, NumPy 2.5.2, SciPy 1.18.1,
scikit-learn 1.9.0, DoubleML 0.11.4 and Tigramite 5.2.10.1. These are observed
execution identities. Only the separately preserved release records establish
the two causal packages' availability before the literature cutoff; runtime
versions of other dependencies are not silently treated as cutoff releases.

\subsection{Deriving the execution choices}
For $N$ independent trajectories, $H$ decisions and $K$ finite states, one
forward/backward pass with dense transitions costs $O(NHK^2)$. With $A$
actions and $M$ signal categories, the parameter count is
\[
 p=(K-1)+K(M-1)+K(A-1)+AK(K-1)+AK.
\]
The current reference retains adjacent-state posterior arrays, requiring
$O(NHK^2)$ memory as well as forward/backward arrays. Batching whole
trajectories can reduce that memory because their likelihoods and scores
add at a common parameter. It cannot replace the shared parameter by
independently optimized shard parameters. Finite-differencing an analytic
score for information adds $2p$ score evaluations and an $O(p^3)$ dense
factorization. Extra starts and $B$ full refits multiply the fitting cost;
the elapsed-time table includes only its stated $B=3$ cost experiment.

Independent trajectories and independent bootstrap fits are natural
parallel units. For a fixed parameter, deterministic reductions can combine
batch log likelihoods, gradients and information, with roundoff checked
against the serial reference. Distributed optimization additionally needs
synchronized parameters and a stopping rule on the global score. Emission
validation terms must enter exactly once, with their original sample weight;
duplicating them on every worker changes the likelihood and understates
calibration uncertainty. A trajectory cannot be cut into independent
fragments without transferring the correct filtering distribution and its
parameter sensitivities across the boundary.

For a response vector with $d$ coordinates, direct dense Bartlett covariance
cost is $O(nLd^2)$ and storage includes the $n\times d$ scores. Prefiltering
adds a $d$-variable regression and recolouring solves. A date partition needs
the original lagged observations, future outcome horizons and score-lag
cross-products that straddle its boundaries. Halo records supply context;
only owned origins contribute to the target sum. Center with the global
mean and retain the cross-boundary terms. Dropping them for faster reduction
changes the covariance. Customer shards under common shocks must combine
their scores by calendar date before a time-series covariance is formed.

Conditional-independence search has a different obstacle. With $v$ variables
and lag window $L$, each node can have order $vL$ candidate parents. Testing
all conditioning sets of size $s$ involves up to
$\binom{vL}{s}$ sets per candidate relation before pruning; a dense
partial-regression calculation also depends on sample size and $s$.
Actual parent selection can remove much of this work, explaining why two
samples with the same dimensions need different time. A search cap is an
algorithmic restriction that can remove oracle completeness, not a purely
physical optimization. Parallel tests must preserve stage dependencies and
orientation rules. Distributed throughput cannot be inferred by multiplying
the one-thread table by the host's 256 logical CPUs.

The resulting architecture uses whole-customer batches for history and
controlled fits, date-aware reductions for common-shock responses, and
explicit stage scheduling for discovery. Every physical plan must reproduce
the declared serial target and numerical tolerance before claiming a speed
benefit. The measured references identify dominant costs and feasible
parallel units. Large distributed storage, fault recovery, accelerator
qualification and institutional workload service levels remain concrete
construction tasks rather than conclusions of these local measurements.


\section{Partition the experiment without changing its information}
\label{s6:partitioned}
A serial reference is especially useful when calibration, nuisance fitting
and dependent inference interact. In the dynamic funding construction,
workers first accumulate independent audit successes and counts by stratum.
The pilot fitting pass accumulates $X'X$, $X'Y$, response sums and squared
sums on the same fixed training half. Validation losses use the same
candidate penalties and pilot holdout. After choosing the penalty, a second
pilot pass forms the full normal equations. The evaluation pass accumulates
the candidate score numerators and denominators by stratum. No evaluation
outcome chooses the learner or covariance bandwidth.

Let $\hat\psi_t$ be the stacked candidate-response influence, and let
$\bar\psi$ be its global mean. For lag $k$, partition the later index $t$:
\[
 \hat\Gamma_k=T^{-1}\sum_b\sum_{t\in I_b,\,t>k}
   (\hat\psi_t-\bar\psi)(\hat\psi_{t-k}-\bar\psi)'.
\]
Each partition receives a left halo of at most $L$ dates, so cross-boundary
products occur once. Computing a covariance inside each partition and then
averaging would omit those products and could center different experiments.
The resulting Bartlett matrix, inverse-linkage Jacobian, audit covariance
and financial weights are applied in the same order as the serial derivation.
Partition counts are not independent-cluster counts.

Eighteen fresh-process runs compare one, two and four workers at 720,
5,760 and 23,040 evaluation origins, with four times as many pilot origins,
48 branches, seven horizons, two candidate portfolios and two audit strata.
Each worker's numerical libraries are limited to one observed thread.
The measured workflow includes synthetic source generation, revision-aware
materialization, calibration, pilot selection/fitting, response fitting,
covariance and financial result construction. The direct reference comparison,
consumer formatting, semantic checks and receipt serialization are excluded
from the reported workflow time.
All serial/partitioned financial and covariance outputs agree with the
separate reference to below $5\times10^{-10}$ in the recorded absolute
comparisons. Floating-point reduction order can change last bits; it cannot
change the data splits or the number of independent dates.

Mean times over two fresh processes are about .195/.282/.367 seconds at
720 origins and 4.818/5.038/5.024 seconds at 23,040 origins, in
one/two/four-worker order. The largest parent process peaks around 474 MiB;
this whole-process peak also includes the untimed reference comparisons and
semantic checks, rather than isolating the timed workflow's memory.
Child peak-RSS counters are also retained, but are process-accounting maxima,
not a measured simultaneous sum across hosts or workers. Parallel execution
is slower here: process and transfer overhead dominate this small fixed
basis. The evidence supports serial execution as the default for these
dimensions. It does not support a claim of distributed speedup.

The partition construction nevertheless qualifies a harder boundary than
reducing precomputed independent scores: audit uncertainty, pilot model
selection, complete nuisance fits and chronological covariance remain in
one experiment. It is still a single-host, in-memory reference. At scale,
the source materialization and sufficient-statistic passes can be streamed,
but network transfer, object storage, fault recovery and concurrent study
load remain unmeasured. Longer POMDP histories and graph search have
different dominant costs and cannot inherit this timing result. Capacity
planning must combine these measured dimensions with the earlier latent
and discovery workload profiles.

\section{Streamed reconstruction and restartable score accounting}
\label{c7:restartable}
A physical execution plan must preserve the statistical experiment after
interruption. The selected calendar source is a sorted retrospective feed:
predecision controls arrive by their decision date, outcomes correspond to
the declared horizon, and realized outcomes arrive by a fixed extraction
cutoff. Streaming materialization validates these clocks and exactly one
record per origin, writes numeric columns to a temporary memory-mapped array,
checks the complete source identity, then commits the completed array.
Incomplete output is not advertised as a valid observation view. The dense
reference parses the same feed and must produce identical columns.

The nuisance learner may still copy those columns into its fitting matrices.
Streaming the source is therefore not a claim of a general out-of-core Lasso.
Where the estimator admits sufficient-statistic passes, the score computation
can stream without changing inference. For scalar scores $s_t$, first obtain
the global mean $\bar s$. A second pass accumulates
\begin{align*}
 P_0=\sum_t(s_t-\bar s)^2,\qquad
 P_k=\sum_{t:\,t,t-k\text{ in same block}}
               (s_t-\bar s)(s_{t-k}-\bar s),\\
 \widehat\Omega=\frac{P_0+2\sum_{k=1}^{L}(1-k/(L+1))P_k}{n}.
\end{align*}
Each partition carries its immutable bytes, fold, chronological offset and
contribution identity. A halo of at most $L$ prior scores supplies boundary
lag products within the same block. It resets between blocks, so no artificial
lag pair bridges a deleted temporal gap.

The checkpoint binds data and procedure identities, completed contribution
prefixes, count, global sum, lag products and halo. Atomic rename commits a
coherent state; a content checksum detects corruption of that state. Retrying
a committed contribution verifies its source and skips its numerical update.
A changed source, changed bandwidth, overlapping contribution or reordered
manifest is rejected. The design is an exactly-once contribution calculation
on one host, not a claim about a distributed transactional service.

Failure injection covers interruptions after committed moment and lag passes,
interruptions before rename, duplicate retry attempts, corrupted checkpoints,
changed inputs and changed procedures. Recovered results agree with a separately
computed dense Bartlett reference to $10^{-12}$. Source fixtures reject late
features, wrong horizons, duplicate origins and outcomes beyond the cutoff.
These checks substantiate the local recovery mechanism. Network partitions,
concurrent writers and remote object-store semantics require their own platform
experiment.

The complete local comparison includes reading, fitting, uncertainty, both
consumers and serialization in fresh processes. It matches direct expert
composition to the shared operations under the same unit, provenance and
inference obligations, asserting identical public output bytes. Separate
workloads vary pair counts, learner dates, audited nonlinear trajectories and
finite Gaussian episodes; peak memory and source/output/checkpoint bytes are
recorded alongside elapsed time. First JIT compilation is included in the
nonlinear interval, while imports and generation of the synthetic source are
outside it. These boundaries matter more than a favorable isolated kernel time.

Local wrappers show computational parity for the shared scientific kernels;
restartable score persistence adds work and does not establish a speedup.
The recommendation is a serial local reference with optional checkpointing
when recoverability matters. A real multi-host comparison must retain identical
samples, outputs and numerical tolerances, introduce storage/network faults,
and measure concurrent complete requests. Analyst time, adoption and business
value require repeated matched tasks by independent users, rather than an
extrapolation from this author's implementation effort.

